\subsection{Manhattan Distance Trick}

This section is about $\mathbb Z^2$. Let $\mathcal L((x, y)) = (x+y, x-y)$. We have the distance functions:

\begin{itemize}
    \item Manhattan distance: $M(p, q) = |p.x - q.x| + |p.y - q.y|$
    \item Chebyshev distance: $C(p, q) = \max(|p.x - q.x|, |p.y - q.y|)$
\end{itemize}

Then:

\begin{itemize}
    \item $\mathcal L(\mathbb Z^2)$ scales $\mathbb Z^2$ by $\sqrt 2$;
    \item $\mathcal L(\mathbb Z^2)$ rotates $\mathbb Z^2$ by $\frac \pi 4$ clock-wise;
    \item For some $p \in \mathbb Z^2$, $\mathcal L(\{q : M(q, p) \leq d\})$ forms an axis-aligned square in $\mathcal L(\mathbb Z^2)$, with bottom-left corner $\mathcal L(p) - (d, d)$ and upper-right corner $\mathcal L(p) + (d, d)$;
    \item $M(p, q) = C(\mathcal L(p), \mathcal L(q))$, and $C(p, q) = M(\mathcal L^{-1}(p), \mathcal L^{-1}(q))$.
\end{itemize}

\subsubsection{Higher Dimensions}

On $\mathbb Z^d$, the last fact generalizes as a $2^{d-1}$ dimension transformation:

$$ \mathcal L(p)_{i} = p_0 + \sum_{j=1}^{d-1} (-1)^{(i>>j\&1)} p_j, $$

that is, the sign of the first coordinate is positive, and the others iterate through all $2^{d-1}$ possibilities.
